% Argument-first working draft, 2026-09-05. Hand-edited.
% Scope: supplementary.md; research/{DEFINITIONS,METHODS,METRICS}.md;
% analyses/013-2026-09-04-matched-intervention-manuscript/README.md.
% Empirical conclusions and contribution claims await the results stage.
\section{Introduction}
\label{sec:introduction}

Activation sparsity offers a route to reducing the cost of Transformer inference.
For a linear transformation $y=xW$, a zero coordinate of $x$ nullifies the
corresponding products, which a specialized kernel can skip. Because activations
depend on the input, the available savings can vary across tokens, layers, and
model components. Training-free methods such as TEAL show that pretrained
language models can tolerate setting small-magnitude activations to zero and
that this sparsity can yield measured decoding speedups~\citep{liu2024teal}.
This establishes a useful starting point: a learned representation already
contains some capacity for sparse execution. Intervening during pretraining
raises a further question: can the learning process shape that capacity to
achieve greater sparsity at the same model quality?

Several approaches establish that adaptation to sparsity is possible.
ProSparse combines ReLU activations, progressive $\ell_1$ regularization, and
threshold shifting during continued training~\citep{song2024prosparse}.
\citet{cetin2026sparser} induce highly sparse FFN activations during training
and develop kernels that exploit them. Q-Sparse extends training-time
sparsification to linear projections in both attention and FFNs through
top-$k$ selection~\citep{wang2024qsparse}. These results motivate studying the
ingredients of sparsification separately. A successful combination establishes
an operating point, but explaining that point requires knowing which ingredients
contribute and how their effects depend on one another.

We study activation sparsification as three separable interventions:
\emph{pressure}, which encourages selected activations toward zero;
\emph{thresholding}, which determines which values become exactly zero; and
\emph{site placement}, which determines where each intervention acts.
Our central question is: \emph{how does the effect of activation pressure on
the quality--sparsity trade-off depend on the gate nonlinearity and the sites
that receive pressure?} Separability here is experimental: pressure targets
can be changed while gates remain fixed, and gates can be changed without
adding pressure. Their effects need not be additive. In particular, extending
the gate set and extending the pressure objective are distinct changes, even
when a practical recipe applies both at the same sites.

A candidate explanation follows from the different roles of these interventions.
Pressure reshapes activation magnitudes; the gate converts some of that change
into exact zeros; and the downstream operation determines the computational
value of those zeros. Moving an activation toward zero need not remove a
product until it crosses a gate's threshold. Moreover, a one-sided gate removes
negative values, whereas a symmetric magnitude gate preserves sufficiently
large values of either sign. The same pressure can therefore interact with
different gates in different ways. Its quality cost also depends on how the
auxiliary update interacts with task learning. This account motivates tracing
an intervention from optimization, through sitewise activation changes, to
operation-level sparsity, and testing whether that account predicts behavior
outside the setting in which it was developed.

To connect local behavior to the complete forward pass, we use
$R_{\mathrm{model}}$, the fraction of declared model-wide matrix-multiplication
products with an exactly zero activation operand. The count includes the
Transformer projections and causally valid attention products, with the dense
language-model head retained in the denominator. Throughout, \emph{model-wide
sparsity} refers to this product-weighted quantity; local activation zero
fractions are reported separately. We also define
$R_{\mathrm{model}}^{\max}$, the analytic all-zero reach ceiling of a selected
site topology for a specified architecture and sequence length. It describes
the work reachable through those sites, independently of learned activation
values. This separates the opportunity offered by a topology from the sparsity
realized by a trained model. Neither quantity is a wall-clock measurement.

We organize the study around a \emph{sparsification intervention ladder}
(Figure~\ref{fig:intervention-ladder}), progressing from a stock GeLU model
through FFN-local ReLU and pressure to gates at branch inputs, the attention
context, and query, key, and value operands. The ladder provides a map of
interventions; attribution comes from matched contrasts within it. Our
experimental setting uses randomly initialized Pythia-family architectures
at 14M, 70M, and 410M parameters on MiniPile~\citep{biderman2023pythia,kaddour2023minipile},
with initialization, data order, and training budget matched within each
comparison. The detailed small-model contrasts and selected larger-model
extensions serve different roles: the former examine individual interventions,
while the latter probe their transfer under a fixed data budget. Model size
also changes the operation weights in $R_{\mathrm{model}}$, so a larger sparsity
fraction alone cannot establish a stronger learned response.

The paper is organized around three empirical questions. First, when does
adding pressure improve the quality--sparsity trade-off, and how does that
effect change with gate and pressure placement? Second, can activation and
optimization diagnostics explain the resulting operation-level changes?
Third, which parts of that explanation transfer across model sizes, and under
what execution conditions can the resulting zeros be exploited? Together,
these questions make the interaction, its explanation, and its limits the
objects of study.
